\documentclass[10pt]{article}
\usepackage[a4paper,bottom = 0.6in,left = 0.75in,right = 0.75in,top = 1cm]{geometry}
\usepackage{graphicx}
\usepackage{amsmath}
\usepackage{array}
\usepackage{enumitem}
\usepackage{wrapfig}
\usepackage{microtype}
\usepackage{titlesec}
\usepackage{graphicx}
\usepackage{booktabs}
\usepackage{url}
\usepackage{enumitem}
\usepackage{palatino}
\usepackage{tabularx}
\fontfamily{SansSerif}
\usepackage{fontawesome5}
\usepackage[hidelinks]{hyperref}
\usepackage{textcomp}

\usepackage{verbatim}
\usepackage{makecell}
\usepackage{pbox}

%for color
\usepackage[usenames, dvipsnames]{color}
\definecolor{myblue}{RGB}{10,0,254}

% Adjusting the margins
\enlargethispage{2\baselineskip} % you can adjust the number based on your content's fit in the page

\newcommand{\xfilll}[2][1ex]{
\dimen0=#2\advance\dimen0 by #1
\leaders\hrule height \dimen0 depth -#1\hfill}
\titleformat{\section}{\large\scshape\raggedright}{}{0em}{}
\renewcommand\labelitemi{\raisebox{0.4ex}{\tiny$\bullet$}}
\renewcommand{\labelitemii}{$\cdot$}
\pagenumbering{gobble}


\usepackage[T1]{fontenc}
\usepackage
%[ansinew]
[utf8]
{inputenc}

\usepackage{color}
\definecolor{mygrey}{gray}{0.99}
\textheight=10in
\raggedbottom

\setlength{\tabcolsep}{0in}
\newcommand{\isep}{-2 pt}
\newcommand{\lsep}{-0.5cm}
\newcommand{\psep}{-0.6cm}
\renewcommand{\labelitemii}{$\circ$}

\pagestyle{empty}
%-----------------------------------------------------------
%Custom commands
\newcommand{\resitem}[1]{\item[{\color[RGB]{10,0,254}$\bullet$}] #1 \vspace{-2pt}}
\newcommand{\resheading}[1]{{\small \colorbox{mygrey}{\begin{minipage}{0.99\textwidth}{\textbf{#1 \vphantom{p\^{E}}}}\end{minipage}}}}
\newcommand{\ressubheading}[3]{
\begin{tabular*}{8.62in}{l @{\extracolsep{\fill}} r}
	\textsc{{\textbf{#1}}} & \textsc{\textit{[#2]}} \\
\end{tabular*}\vspace{-8pt}}


\begin{document}

\begin{table}
    \begin{center}
    {\Huge \scshape Oleksandr Kovalyk Borodyak} \\ \vspace{5pt}
    TELEMATICS ENGINEER \\ \vspace{5pt}
	% \small \raisebox{-0.1\height}\faPhone\ +34 645043899 
	\href{mailto:oleksandr.kovalyk@edu.upct.es}{\raisebox{-0.2\height}\faEnvelope\  \underline{oleksandr.kovalyk@edu.upct.es}} ~ 
    \href{https://www.linkedin.com/in/oleksnadr-kovalyk-4a5861213}{\raisebox{-0.2\height}\faLinkedin\ \underline{Oleksandr Kovalyk}}  ~
    % \href{https://github.com/Kovaxs}{\raisebox{-0.2\height}\faGithub\ \underline{Kovaxs}} ~
    % \href{https://websiteurl.com}{\raisebox{-0.2\height}\faGlobe\ \underline{/websiteurl.com}}
    \href{https://orcid.org/0000-0003-4815-8248}{\raisebox{-0.2\height}\faOrcid\ \underline{0000-0003-4815-8248}}
    \vspace{-8pt}
\end{center}
\end{table}    


\setlength{\tabcolsep}{11pt} % decrease the value if the content is overfitting
\begin{table}
\centering
\begin{tabular}{lllll}
\toprule
\textbf{Title}   & \textbf{Institute}    & \textbf{Year} \\ 
\toprule
Master's Degree in Telematics Engineering & Technical University of Cartagena & 2021  \\
Bachelor Degree in Telematics Engineering & Technical University of Cartagena & 2020  \\
\bottomrule \\[-0.75cm]
\end{tabular}
\end{table}
\vspace{-8pt}

%%% SHORT SUMMARY %%%
\section*{{ \LARGE \color{myblue} Summary}\xfilll[0pt]{0.5pt}}
\vspace{-15pt}
\vspace{2.5mm}
% Describe your professional self in a few lines
I am passionate about technology and innovation.
I see myself as a responsible individual, capable of adapting and learning swiftly.
I enjoy working in a team and contributing new ideas.
I consider myself to be a person with initiative, always eager to learn and improve day by day.
I love tinkering with technology and learning new things beyond my field of study. 
I am a person who is always looking for new challenges and opportunities to grow professionally.
In my free time, I usually do sport, my favorite sport is weightlifting.
\vspace{-8pt}

%%%% PROFESSIONAL EXPERIENCE %%%%
\section*{{\LARGE \color{myblue}Professional Experience}\xfilll[0pt]{0.5pt}}
\vspace{-15pt}
\vspace{3.5mm}

%\begin{itemize}[itemsep = -0.75 mm, leftmargin=*]
\textbf{\large Research Assistant}| \normalfont{Technical University of Cartagena \& EU Life Project} \hfill 
%\\[-0.5cm]
% \emph{Awarded Letter of Recommendation for excellent performance during the internship period } 
%\\[-0.5cm]	
	\begin{itemize}[itemsep = -0.75 mm, leftmargin=*]
	\item[{\color[RGB]{10,0,254}$\bullet$}] \noindent Reducing the impact of underwater noise on the marine environment of the Port of Cartagena.\hfill {\href{https://webgate.ec.europa.eu/life/publicWebsite/project/LIFE20-ENV-ES-000387/reducing-the-impact-of-underwater-noise-on-the-marine-environment-of-the-port-of-cartagena}{\small \underline{Project Link} }} 
	\end{itemize} 
\textbf{\large Research Assistant}| \normalfont{Technical University of Cartagena \& MTorres} \hfill { \small February 2022 - May 2023}
%\\[-0.5cm]
 %\emph{Awarded Letter of Recommendation for excellent performance during the internship period } 

	\begin{itemize}[itemsep = -0.75 mm, leftmargin=*]
	\item[{\color[RGB]{10,0,254}$\bullet$}] \noindent Detection of carbon fiber wrapping failures using images from a profilometer.
	\end{itemize} 
\textbf{\large Research Assistant}| \normalfont{Technical University of Cartagena \& \\ \hspace*{3.54cm} | Hospital General Universitario Reina Sofía} \hfill { \small December 2019 - December 2020} 
%\\[-0.5cm]
% \emph{Awarded Letter of Recommendation for excellent performance during the internship period } 
%\\[-0.5cm]	
	\begin{itemize}[itemsep = -0.75 mm, leftmargin=*]
	\item[{\color[RGB]{10,0,254}$\bullet$}] \noindent Automatic feature extraction of the optical nerve of both eyes by
	image processing and statistical pattern recognition for its integration in a telemedicine CAD (Computer-Aided Diagnosis) application. 
	\item[{\color[RGB]{10,0,254}$\bullet$}] \noindent This work has been partially funded by Spanish National projects AES2017-PI17/00771 and
	AES2017-PI17/00821 (Instituto de Salud Carlos III) and regional project 20901/PI/18 (Fundación Séneca). 
	\end{itemize} 


\vspace{-15pt}

%%% SKILLS %%%
\section*{{\LARGE \color{myblue}Technical Skills}\xfilll[0pt]{0.5pt}}

\vspace{-20pt}
\setlength{\tabcolsep}{9pt}
\begin{table}[h]

\begin{tabular}{lllll}

\textbf{Programming} & Python, SQL \\[0.5mm]
\textbf{Development} & XGBoost, Scikit-learn, Pandas, Numpy, Matplotlib, Seaborn, Keras, Tensorflow, Pytorch \\[0.5mm] 
\textbf{Tools} & Git, GitHub, PostgreSQL, Minio, MLflow, Neovim \\[0.5mm]
\textbf{Other} & Docker, Kubernetes, Linux, Proxmox, Latex 

\end{tabular}
\end{table}
% \fi 

\vspace{-20pt}
%%% PROJECTS %%%
\section*{{ \LARGE \color{myblue} Projects Involvement}\xfilll[0pt]{0.5pt}}
\vspace{-13pt}
\vspace{3.5mm}

{\LARGE \textcolor{black}{Software Development}}
\vspace{-1.5pt}
\par\noindent\hrulefill % Add a horizontal line
\\[0.35cm]
\textbf{Reducing the impact of underwater noise on the marine environment of the Port of Cartagena.} \hfill  
\\[0.15cm] 
\normalfont{Python, Docker, Docker-Compose, Minio, PostgreSQL} \hfill { \small March 2022 - Present}
\\[0.15cm]
\emph{Technical University of Cartagena \& EU Life Project} \hfill    
%\\[-0.6cm]
	\begin{itemize}[itemsep = -0.75 mm, leftmargin=*]
	\item[{\color[RGB]{10,0,254}$\bullet$}] \noindent The project consists on build a system based on microservices architecture that capture
	the date from diverse sources and stores it in a database backend.
	\item[{\color[RGB]{10,0,254}$\bullet$}] \noindent The system will be deployed in as a Docker-Compose and consist from a following modules:
	\begin{itemize}
		\item \textbf{Data retrieving}: module that retrieves data from two hydrophones and AIS system and store it in a database.
		\item \textbf{Data processing}: module that intake the data and process it to extract the noise level and associate it with the AIS data.
		\item \textbf{API REST}: module that provides an API to access the data.
	\end{itemize}
	\end{itemize} 	

\vspace{0.5cm}
\noindent{\LARGE \textcolor{black}{Machine Learning}}
\vspace{-1.5pt}
\par\noindent\hrulefill % Add a horizontal line
\\[0.35cm]
\noindent\textbf{Detection of carbon fiber wrapping failures using images from a profilometer.} 
\\[0.15cm]
\normalfont{Python, Keras, Tensorflow, RetinaNet} \hfill {\ \small Feb 2022 - May 2023}
\\[0.15cm]
\emph{Technical University of Cartagena \& MTorres} \hfill 
%\\[-0.6cm]
	\begin{itemize}[itemsep = -0.75 mm, leftmargin=*]
	\item[{\color[RGB]{10,0,254}$\bullet$}] \noindent 	The project involves building a model and its corresponding pipeline to detect carbon
	fiber wrapping failures using images from a profilometer. The project have the following stages:
	\begin{itemize}
		\item Data analysis and preprocessing pipeline.
		\item Selecting of the performance metrics.
		\item Train RetinaNet model implemented in Python with Keras (Tensorflow backend) and hyperparameter tuning.
		\item Model evaluation and final results. 
	\end{itemize}
	\end{itemize} 	

\noindent\textbf{Automatic extraction of optic nerve features from both eyes through image processing and recognition.}
\\[0.15cm]
\normalfont{Python, Keras, Tensorflow} \hfill {\ \small Dec 2019 - Dec 2020}
\\[0.15cm]
\emph{Technical University of Cartagena \& Hospital General Universitario Reina Sofía} \hfill  
%\\[-0.6cm]
	\begin{itemize}[itemsep = -0.75 mm, leftmargin=*]
	\item[{\color[RGB]{10,0,254}$\bullet$}] \noindent In this project my task was to curate the dataset obtained from the Hospital General Universitario Reina Sofía and develope a DL-based APP
	that take two ocular fundus images to diagnose the glaucoma.
	\item[{\color[RGB]{10,0,254}$\bullet$}] \noindent The project is part of Automatic feature extraction of the optical nerve of both eyes by
	image processing and statistical pattern recognition for its integration in a telemedicine CAD (Computer-Aided Diagnosis) application. 
	\end{itemize} 	

\vspace{0.5cm}
\noindent{\LARGE \textcolor{black}{Personal Projects}}
\vspace{-1.5pt}
\par\noindent\hrulefill % Add a horizontal line
\\[0.35cm]
\noindent\textbf{Kubernetes cluster building on Proxmox and managed with Ansible.} 
\\[0.15cm]
\normalfont{Ansible, Kubernetes, Proxmox, Docker} \hfill 
\\[0.15cm]
\emph{Personal} \hfill 
%\\[-0.6cm]
	\begin{itemize}[itemsep = -0.75 mm, leftmargin=*]
	\item[{\color[RGB]{10,0,254}$\bullet$}] \noindent In this project, my goal is to replicate a real-word scenario of deploying a Machine Learning model within a Kubernetes cluster.
	For this purpose, I have built a Kubernetes cluster on Proxmox, managing the VMs with Ansible. I am also use the Kubernetes cluster to learn new technologies such as Kafka, Spark, DeltaLake, etc.
	The cluster serve as a playground for me to learn new technologies and experiment with them.
	\end{itemize} 	

\vspace{-15pt}

%%% PUBLICATIONS %%%
\section*{{\LARGE \color{myblue}Publications}\xfilll[0pt]{0.5pt}}
\vspace{-8pt}

\noindent\textbf{\large{Decision trees for glaucoma screening based on the asymmetry of the retinal nerve fiber layer in optical coherence tomography}} \hfill {\ \small Jun 2022}
\\[0.15cm]
\emph{MDPI} \hfill{\small DOI: 10.3390/s22134842}
%\\[-0.6cm]
\begin{itemize}[itemsep = -0.75 mm, leftmargin=*]
    \item[{\color[RGB]{10,0,254}$\bullet$}] \noindent \textbf{Authors:} Berenguer-Vidal, Rafael; Verd\'u-Monedor, Rafael; Morales-S\'anchez, Juan; Sell\'es-Navarro, Inmaculada;
	Kovalyk Oleksandr; Sancho-G\'omez, Jos\'e-Luis.
\end{itemize}

\noindent\textbf{\large{Improving glaucoma diagnosis assembling deep networks and voting schemes}} \hfill {\ \small Jun 2022}
\\[0.15cm]
\emph{MDPI} \hfill{\small DOI: 10.3390/diagnostics12061382}
%\\[-0.6cm]
\begin{itemize}[itemsep = -0.75 mm, leftmargin=*]
    \item[{\color[RGB]{10,0,254}$\bullet$}] \noindent \textbf{Authors:} S\'anchez-Morales, Adrian; Morales-S\'anchez, Juan; Kovalyk, Oleksandr; Verd\'u-Monedero, Rafael;
	Sancho-G\'omez, Jos\'e-Luis.  
\end{itemize}

\noindent\textbf{\large{PAPILA: Dataset with fundus images and clinical data of both eyes of the same patient for glaucoma assessment}} \hfill {\ \small Jun 2022}
\\[0.15cm]
\emph{Nature} \hfill{\small DOI: 10.1038/s41597-022-01388-1}
%\\[-0.6cm]
\begin{itemize}[itemsep = -0.75 mm, leftmargin=*]
    \item[{\color[RGB]{10,0,254}$\bullet$}] \noindent \textbf{Authors:} Kovalyk, Oleksandr; Morales-S\'anchez, Juan; Verd\'u-Monedero, Rafael; Sell\'es-Navarro, Inmaculada;
	Palaz\'on-Cabanes, Ana; Sancho-G\'omez, Jos\'e-Luis.
\end{itemize}

%%% CONFERENCES %%%

\section*{{\LARGE \color{myblue}Conferences}\xfilll[0pt]{0.5pt}}
\vspace{-13pt}
\vspace{3.5mm}

\noindent{\LARGE \textcolor{black}{International Conferences}}
\vspace{-1.5pt}
\par\noindent\hrulefill % Add a horizontal line

\noindent\textbf{\large{Analysis of the asymmetry in RNLF thinkness using sprectralis OCT measurements in healthy and glaucoma patients}} \hfill {\ \small May 2022}
\\[0.15cm]
\emph{Springer} \hfill{\small DOI: 10.1007/978-3-031-06242-1\_50}
%\\[-0.6cm]
\begin{itemize}[itemsep = -0.75 mm, leftmargin=*]
	\item[{\color[RGB]{10,0,254}$\bullet$}] \noindent \textbf{Name of the conference:} 9th International Work-Conference on the Interplay Between Natural and Artificial Computation
	\item[{\color[RGB]{10,0,254}$\bullet$}] \noindent \textbf{Location of the event:} Puerto de la Cruz, Tenerife, Spain
	\item[{\color[RGB]{10,0,254}$\bullet$}] \textbf{Date of the event:} May 31 - June 3, 2022
    \item[{\color[RGB]{10,0,254}$\bullet$}] \noindent \textbf{Authors:} Berenguer-Vidal, Rafael; Verd\'u-Monedero, Rafael; Morales-S\'anchez, Juan; Sell\'es-Navarro,
	Inmaculada; Kovalyk Oleksandr.
\end{itemize}

\noindent{\LARGE \textcolor{black}{National Conferences}}
\vspace{-1.5pt}
\par\noindent\hrulefill % Add a horizontal line

\noindent\textbf{\large{Automatic Instance Segmentation Applied to Retinography for Optic Disc and Cup Segmentation}} 
\begin{itemize}[itemsep = -0.75 mm, leftmargin=*]
	\item[{\color[RGB]{10,0,254}$\bullet$}] \noindent \textbf{Name of the conference:} XXXVIII Annual Congress of the Spanish Society of Biomedical Engineering
	\item[{\color[RGB]{10,0,254}$\bullet$}] \noindent \textbf{Location of the event:} Valladolid, Castilla y Le\'on, Spain	
	\item[{\color[RGB]{10,0,254}$\bullet$}] \noindent \textbf{Date of the event:} 25-27 November 2020
	\item[{\color[RGB]{10,0,254}$\bullet$}] \noindent \textbf{Organizing entity:} University of Valladolid 
	\item[{\color[RGB]{10,0,254}$\bullet$}] \noindent \textbf{Authors:} Kovalyk, Oleksandr; Morales-S\'anchez, Juan; Verd\'u-Monedero, Rafael; Sell\'es-Navarro, Inmaculada. 
\end{itemize}

\noindent\textbf{\large{Application of mobileNet for early diagnosis of glaucoma: a binocular approach}} 
\begin{itemize}[itemsep = -0.75 mm, leftmargin=*]
	\item[{\color[RGB]{10,0,254}$\bullet$}] \noindent \textbf{Name of the conference:} XLI Annual Congress of the Spanish Society of Biomedical Engineering
	\item[{\color[RGB]{10,0,254}$\bullet$}] \noindent \textbf{Location of the event:} Cartagena, Murcia, Spain	
	\item[{\color[RGB]{10,0,254}$\bullet$}] \noindent \textbf{Date of the event:} 22-24 November 2023
	\item[{\color[RGB]{10,0,254}$\bullet$}] \noindent \textbf{Organizing entity:} Technical University of Cartagena 
	\item[{\color[RGB]{10,0,254}$\bullet$}] \noindent \textbf{Authors:} Kovalyk, Oleksandr; Morales-S\'anchez, Juan; Verd\'u-Monedero, Rafael; Sell\'es-Navarro, Inmaculada;
	Sancho-G\'omez, Jos\'e-Luis. 
\end{itemize}

%%% COURSES %%%

% \section*{{\LARGE \color{myblue}Key Courses Undertaken}\xfilll[0pt]{0.5pt}}
% \vspace{-15pt}
% \vspace{1.5mm}
% 
% \textbf{\large Mathematics and Statistics}
% \\[-0.6cm]
% \begin{itemize}[itemsep = -0.75 mm, leftmargin=*]
% 	\item[{\color[RGB]{10,0,254}$\bullet$}] \noindent Calculus, Probability and Random Processes, Basics of Statistics, Data Preprocessing for Regression Analysis,
% Regression Models, Classification Models, Support Vector Machines, Ensemble Techniques, Artificial Neural Networks
% 	\end{itemize}
% \textbf{\large Computer Science}
% \\[-0.6cm]
% \begin{itemize}[itemsep = -0.75 mm, leftmargin=*]
% 	\item[{\color[RGB]{10,0,254}$\bullet$}] \noindent Structured Programming Approach, Data Structures and Algorithms, Operating Systems, Computer Networks
% 	\end{itemize}
% \vspace{-15pt}


%%% POSITIONS OF RESPONSIBILITY %%%
%  \section*{{\LARGE \color{myblue}Positions of Responsibility}\xfilll[0pt]{0.5pt}}
% \vspace{-15pt}
% \vspace{1.5mm}
% 
% \textbf{\large{Position}}} \hfill {\ \small Feb 2022 - Present}
% \\[0.05cm]
% \emph{Name of Organization}
% \\[-0.6cm]
% \begin{itemize}[itemsep = -0.75 mm, leftmargin=*]
%     \item[{\color[RGB]{10,0,254}$\bullet$}] \noindent Your item 1.
% 	\item[{\color[RGB]{10,0,254}$\bullet$}] \noindent Your item 2.
% 	\end{itemize}
% \textbf{\large{Position}} \hfill {\ \small Feb 2022 - Present}
% \\[0.05cm]
% \emph{Name of Organization}
% \\[-0.6cm]
% \begin{itemize}[itemsep = -0.75 mm, leftmargin=*]
%     \item[{\color[RGB]{10,0,254}$\bullet$}] \noindent Your item 1.
% 	\item[{\color[RGB]{10,0,254}$\bullet$}] \noindent Your item 2.
% 	\end{itemize}
%  \textbf{\large{Position}} \hfill {\ \small Feb 2022 - Present}
% \\[0.05cm]
% \emph{Name of Organization}
% \\[-0.6cm]
% \begin{itemize}[itemsep = -0.75 mm, leftmargin=*]
%     \item[{\color[RGB]{10,0,254}$\bullet$}] \noindent Your item 1.
% 	\item[{\color[RGB]{10,0,254}$\bullet$}] \noindent Your item 2.
% 	\end{itemize}
% 

%%% PUBLICATIONS %%%
% \section*{{\LARGE \color{myblue}Publications}\xfilll[0pt]{0.5pt}}
% \vspace{-8pt}
% \textbf{\large{Title of Publication}} \hfill {\ \small Mar 2022}
% \\[0.05cm]
% \emph{Name of Publishing Organization} \hfill{\small ISBN: 123-456-789-0}
% \\[-0.6cm]
% \begin{itemize}[itemsep = -0.75 mm, leftmargin=*]
%     \item[{\color[RGB]{10,0,254}$\bullet$}] \noindent Your item 1.
%     \item[{\color[RGB]{10,0,254}$\bullet$}] \noindent Your item 2.
% \end{itemize}
% 
% 
% %%% ACHIEVEMENTS %%%
% \section*{{\LARGE \color{myblue}Achievements \& Certifications}\xfilll[0pt]{0.5pt}}
% \vspace{-15pt}
% \vspace{1.5mm}
% 
% \begin{itemize}[itemsep = -0.75 mm, leftmargin=*]
% \item[{\color[RGB]{10,0,254}$\bullet$}] \noindent Achievement 1.
% \item[{\color[RGB]{10,0,254}$\bullet$}] \noindent Achievement 2.
% \item[{\color[RGB]{10,0,254}$\bullet$}] \noindent Achievement 3.
% \end{itemize}

\end{document}
